\section*{Ethical Considerations}

This paper studies an attack on safety-critical learning-enabled control.  All
experiments run in simulation using Safe-Control-Gym with PyBullet dynamics
and a public third-party gridworld implementation.  No physical plant,
vehicle, deployed industrial system, human subject, or personal data was
involved.  Reported failures are simulated constraint violations used to
compare deployment contracts.

The evaluated threat begins \emph{after} compromise of a reward-ingestion
principal in the asynchronous learning-data plane and grants only bounded
scalar writes to reward records.  Initial access, product-specific
exploitation, and evasion of log-integrity controls fall outside the artifact
and evaluation.  The experiments characterize the consequence of a
least-privilege data-integrity gap for evidence inheritance at raw-policy
release.  We therefore describe a composition weakness in an assurance
workflow rather than an exploit against an identified fielded product, and we
publish no capability that lowers the cost of obtaining the assumed access.

Stakeholders are operators of runtime-assured learning controllers, vendors of
shielding and safety-filter tooling, and the public exposed to the controlled
plant.  The main foreseeable harm from publication is that an attacker who
already holds reward-ingestion access learns that a retired shield leaves
latent policy damage exposed.  We judge the disclosure benefit to dominate:
the exposure follows from documented assurance architectures, the attack is
detected by the simple task-aware reward checks we measure, and the paper
states deployable countermeasures at both ends of the authority transition.
Origin-bound reward records and trusted recomputation remove the upstream
precondition; resident predictive authority, commit admission, and permanent
filtering contain the downstream physical consequence.

The third-party lifecycle we instrument is a published research artifact
released under an open license by its authors, used unmodified except for
documented instrumentation and reward-record patches that the artifact
itemizes.  No vendor product, hosted service, or operational deployment was
tested, so no coordinated vulnerability disclosure was owed to a third party.
We follow the USENIX Security ethics guidelines and the principles of the
Menlo Report in reporting the above.

\section*{Open Science}

We support the USENIX Security open science policy and will make every
artifact required to reproduce the reported results openly available.  The
artifact contains: (i) the experiment drivers for the controlled
Safe-Control-Gym study and the third-party shield-retirement case study;
(ii) row-level result files backing every table and figure, with source-linked
provenance columns; (iii) exact commands, run namespaces, upstream revisions,
model files, and dependency versions; (iv) locked protocol documents recording
the analysis chronology; and (v) SHA-256 checksums for all locked files.
Automated checks in the artifact audit table provenance, claim locks, and
manuscript format.

For anonymous review, the artifact is submitted as an anonymized archive
through the submission system, so the review committees can access every
component without contacting the authors.  Upon acceptance we will publish the
same artifact in a public repository with a permanent DOI under an open
license, together with the third-party upstream revision pins required to
rebuild the case-study environment.  The artifact retains the scope boundaries
stated in the paper: the trusted-anchor, certificate-coverage, timing, PPO,
benign-utility, and protocol-chronology limitations are recorded with the data
rather than removed.
